\documentclass{ca-capstone}
\capstonerefs{knuth-taocp2,wiki-lfsr,sagemath}
\begin{document}
% Generated by `make bench CAPSTONE=lcg-lfsr-cryptanalysis` (taskgen.capstone report) from
% bench/results.json: reference time x slack on the reference machine.
% Do not edit: change bench/config.yaml and re-run the benchmark.
\newcommand{\ReqSlack}{10}
\newcommand{\ReqBitsMax}{32}
\newcommand{\ReqSeconds}{0{,}01}

\capstonetitle{Криптоанализ ЛКГ и РСЛОС}

\goal
Требуется реализовать генераторы псевдослучайных последовательностей, исследовать их периоды
и алгоритмы восстановления параметров и состояния по выходной последовательности.

\functionality
\begin{enumerate}
  \item Линейный конгруэнтный генератор (ЛКГ): генерация, вычисление периода и длины
    предпериода, проверка условий полного периода по заданной факторизации модуля.
  \item Восстановление следующего состояния при известных параметрах и восстановление
    параметров ЛКГ по последовательным полным выходам на однозначно определённых примерах.
  \item РСЛОС над $\GF{2}$: генерация, вычисление периода и длины предпериода, переход между
    рекуррентным уравнением, многочленом связей и характеристическим многочленом.
  \item Алгоритм Берлекэмпа-Месси и восстановление РСЛОС по не менее чем $2L$ битам выхода.
  \item Реализация KPA-атаки на поточный алгоритм шифрования на основе XOR и каждого из генераторов.
  \item Тесты вырожденных рекуррентностей, неполных периодов и необратимых параметров ЛКГ
    в операциях, где они допустимы.
\end{enumerate}

\requirements
\begin{enumerate}
  \item Восстановленные параметры должны независимо воспроизводить наблюдения; линейная
    сложность и характеристический многочлен проверяются внешним оракулом.
  \item Входы для восстановления гарантированно имеют однозначный ответ. Неоднозначные случаи
    можно отклонять; их перечисление и анализ коротких наблюдений не требуются.
  \item В реализации атаки по известному открытому тексту для заданного генератора на вход подаются
    шифротекст и известное начало открытого текста; требуется восстановить параметры использованного
    генератора и расшифровать оставшуюся часть зашифрованного сообщения.
  \item Смешанные запросы с модулем ЛКГ длиной до \textbf{\ReqBitsMax~бит} должны выполняться
    не более чем за \textbf{\ReqSeconds~с}; поиск периодов ограничен отдельно заданными классами
    малых модулей ЛКГ и небольшой длины РСЛОС.
\end{enumerate}

\details
Разрешается использовать любой язык программирования, в частности Python, систему 
компьютерной алгебры SageMath и встроенные большие целые. Библиотечный
Берлекэмп-Месси и готовые средства восстановления генераторов в решении запрещены и
используются только как оракул. Факторизация модуля для условий полного периода входит во
входные данные.

Атаки на усечённые выходы, комбинированные генераторы и корреляционные атаки не входят в
задание. Следует сохранять невязки Берлекэмпа-Месси и промежуточные многочлены.

\schedule
Работа выполняется индивидуально. Срок выполнения -- \textbf{месяц}; ожидаемая трудоёмкость
-- примерно 35-39 часов. Реализация должна быть разработана и написана вами.

\remarks
\begin{enumerate}
  \item Набор тестов не должен содержать скрытых неоднозначных задач восстановления.
  \item Расшифрование проверяется повторным шифрованием с независимо вычисленной гаммой.
  \item Оптимизация кода не требуется.
\end{enumerate}

\references
\end{document}
